\section{Prototype Description}
In the following section an overall description of the prototype is given. This is done to help the reader understand the utilized setup in the project.
The prototype is divided into 3 separate subsystems. The quadcopter, the ground station and the vicon system. These subsystems can be seen in \figref{prototypediagram}. 

\begin{figure}[H] 
	\includegraphics[scale=.5]{figures/prototypediagram}
	\centering
	\captionsetup{justification=centering}
	\captionof{figure}{Schematic diagram of the prototype, which includes the quadcopter, the ground station and the vicon system.}
	\label{prototypediagram}
\end{figure}

%!!explain somewhere what vicon is and that it can be found in aalborg university, maybe we should have a little section, together with all the other hardware?!!

For this project the Vicon system is utilized to capture the position and orientation of the quadcopter. This is to avoid time consuming sensor set up on the quadcopter, as the Vicon system is already configured and good to go. The Vicon system is describe in further detail in section ??. The Vicon system captures the quadcopters position, in form of coordinates, and its orientation, in the form of euler angles. The information is then transmitted to the Ground Station.

The Ground Station utilizes the data it receives from the Vicon system to calculate the linear velocity and linear acceleration of the quadcopter. References commands for controller, located on the quadcopter, is also needed and is generated from the ground station. These are transmitted together with the position, orientation, linear velocity and linear acceleration to the quadcopter. This is transmitted by utilizing a generated network interface and a wireless module.

The quadcopter consist of a microcontroller, four motorcontrollers, four motors, a battery and a wireless module. The microcontroller main tasks is to handle the network between itself and the ground station, the control calculations and to generate the control signals which should be transmitted to the four motorcontrollers. \\To be able to calculate the rotational speed needed of the motor, the five variables received from the ground station is needed, i.e the control reference, position, orientation, linear velocity and linear acceleration of the quadcopter. Hereafter it is possible to calculate the necessary PWM signals needed, by utilizing the rotational speed calculated by the controller and the measured voltage on the battery. \\Futhermore the quadcopter should also be able to send the battery voltage level to the ground station in order to detect when to end the flight.

%The quadcopter's most complex component is the microcontroller. It handles the control calculations, the PWM calculations according to the required control action and battery level, and the communication protocol necessary to exchange data with the ground station.

%The microcontroller requires 4 PWM outputs with different and independent duty cycles for controlling the 4 motors in the quadcopter. 

%The network interface in the microcontroller handles the incoming data from the ground station. This data consist of reference commands for the controllers and sensor data. It should also be able to send the battery voltage level to the ground station in order to detect when to end the flight.

%The communication between the ground station and the quadcopter is done through wireless modules. These should allow the implementation of the communication by designing only the transport layer protocol.

%The ground station gathers the information coming from the Vicon system and sends it to the quadcopter. This also includes doing calculations with the received Vicon data to obtain other variables like velocities or accelerations.

%Finally, Vicon system uses the position of the markers placed in the quadcopter to provide its position and orientation. The orientation data is given in multiple forms but only Euler angles data is considered.


%\textbf{Battery should be moved somewhere else}
%The battery, which is available for the prototype, is a Zippy Flightmax battery. Its weights 141 gram, has an capacity of 1500 mAh, a voltage of 11.1 volts and a discharge current of 20 amperes. 



